\begin{python0}
from solutions import *; clear()
\end{python0}
\sectionthree{Review: C-strings}

You already know that a string is an array of \texttt{char} terminated by '\verb!\0!'. '\verb!\0!' has numeric value 0. Here's an example:

\begin{console}
char x[] = "abc"; // same as {'a', 'b', 'c', '\0'}
\end{console}

These are called C-strings, i.e., C-strings are character arrays
containing the null character.

There are some useful functions to work with C strings of which some are
listed below (see CISS240):

\begin{console}
strcpy(s1, s2) // copy
strcmp(s1, s2) // compare
strcat(s1, s2) // concat
strlen(s1) // string length
\end{console}

Note that in order to work with the above functions, you must include
the header \texttt{cstring} like so:

\begin{console}
#include <cstring>
\end{console}

Note that the relational operators ==, !=, etc. don't work with the string content but rather with the pointer to the first value the array. In particular, note that \texttt{s1 == s2} compares the addresses of \texttt{s1} and \texttt{s2}. In other words, it does not compare the contents of the arrays. Likewise for $!=$, $<$, etc.

